\documentclass[11pt,a4paper]{article}

\usepackage[utf8]{inputenc}
\usepackage[T1]{fontenc}
\usepackage[english]{babel}
\usepackage[a4paper,margin=2.1cm]{geometry}
\usepackage{microtype}
\usepackage{booktabs}
\usepackage{tabularx}
\usepackage{array}
\usepackage{xcolor}
\usepackage{enumitem}
\usepackage{listings}
\usepackage{hyperref}
\usepackage{fancyhdr}
\usepackage{titlesec}

\definecolor{ansblue}{HTML}{1E3A8A}
\definecolor{anslight}{HTML}{F3F6FC}
\definecolor{ansgray}{HTML}{5B6472}
\definecolor{anscode}{HTML}{F8FAFC}
\definecolor{ansgreen}{HTML}{0F766E}
\definecolor{ansred}{HTML}{B91C1C}

\hypersetup{
  colorlinks=true,
  linkcolor=ansblue,
  urlcolor=ansblue,
  pdftitle={ANS Part 1 Exam Prep},
  pdfauthor={Enes Ozkan},
  pdfsubject={CSE 341 Part 1 Exam Preparation},
  pdfkeywords={ANS, CSE341, programming languages, lexer, parser, Sebesta}
}

\setlength{\parindent}{0pt}
\setlength{\parskip}{5pt plus 1pt minus 1pt}
\setlist[itemize]{leftmargin=1.2em,itemsep=2pt,topsep=3pt}
\setlist[enumerate]{leftmargin=1.5em,itemsep=2pt,topsep=3pt}

\titleformat{\section}{\Large\bfseries\sffamily\color{ansblue}}{\thesection}{0.7em}{}
\titleformat{\subsection}{\large\bfseries\sffamily}{\thesubsection}{0.7em}{}
\titleformat{\subsubsection}{\normalsize\bfseries\sffamily}{\thesubsubsection}{0.7em}{}

\pagestyle{fancy}
\fancyhf{}
\fancyfoot[L]{\footnotesize\color{ansgray}CSE 341 ANS Part 1 Exam Prep \(\cdot\) Enes \"Ozkan}
\fancyfoot[R]{\footnotesize\color{ansgray}\thepage}
\renewcommand{\headrulewidth}{0pt}

\newcommand{\code}[1]{\texttt{\detokenize{#1}}}
\newcommand{\important}[1]{\textbf{\color{ansred}#1}}
\newcommand{\design}{\textbf{Design intent}}
\newcommand{\actual}{\textbf{Actual D2 behavior}}

\newcolumntype{Y}{>{\raggedright\arraybackslash}X}

\lstdefinelanguage{ANS}{
  morekeywords={let,function,return,for,in,if,elif,else,print,
    element,compound,reaction,reactants,products,formula,name,
    atomic_mass,valence,symbol,check,balanced,energy,exothermic,endothermic},
  sensitive=true,
  morecomment=[l]{//},
  morestring=[b]",
  morestring=[b]'
}

\lstdefinestyle{ans}{
  language=ANS,
  basicstyle=\ttfamily\small,
  keywordstyle=\bfseries\color{ansblue},
  commentstyle=\itshape\color{ansgray},
  stringstyle=\color{ansgreen},
  backgroundcolor=\color{anscode},
  showstringspaces=false,
  columns=fullflexible,
  keepspaces=true,
  breaklines=true,
  frame=single,
  rulecolor=\color{anslight},
  framerule=0.4pt,
  xleftmargin=4pt,
  xrightmargin=4pt
}

\lstdefinestyle{csmall}{
  language=C,
  basicstyle=\ttfamily\small,
  keywordstyle=\bfseries\color{ansblue},
  commentstyle=\itshape\color{ansgray},
  stringstyle=\color{ansgreen},
  backgroundcolor=\color{anscode},
  showstringspaces=false,
  columns=fullflexible,
  keepspaces=true,
  breaklines=true,
  frame=single,
  rulecolor=\color{anslight},
  framerule=0.4pt,
  xleftmargin=4pt,
  xrightmargin=4pt
}

\begin{document}

\begin{center}
  {\footnotesize\color{ansgray}Gebze Technical University \(\cdot\) CSE 341: Concepts of Programming Languages\par}
  \vspace{8pt}
  {\Huge\bfseries\sffamily ANS Part 1 Exam Prep\par}
  \vspace{4pt}
  {\Large Alchemist's Notation Script\par}
  \vspace{8pt}
  {\normalsize Student: \textbf{220104004931 Enes \"Ozkan}\par}
  {\normalsize Exam 1: \textbf{Thursday, May 14, 2026, 08:30} \(\cdot\) 90 minutes\par}
  \vspace{8pt}
  {\small Prepared from the submitted PDFs, D2 C code, D2/D3 tests, and Sebesta Chapters 1, 3, 4, and 5.}
\end{center}

\vspace{8pt}
\tableofcontents
\newpage

\section{How To Use This Document}

This document is written for you, Enes, as a crash course and defense sheet for your ANS Part 1 exam. The exam can ask about your own submission, so the important skill is not memorizing isolated facts. You need to explain how the language design, grammar, lexer, parser, AST, tests, and Sebesta terminology connect.

\textbf{Main rule for answers:} separate what Part 1 actually implements from what the language design promises for Part 2. Part 1 is a lexer and parser with AST output. It does not execute programs, type-check them, or verify balanced reactions yet.

\subsection{Your Claimed Areas}

The D6 contribution report says your main ownership areas are:

\begin{itemize}
  \item D1 4.2 Lexical Structure.
  \item D1 4.6 Names, Binding, Scope, and Lifetime.
  \item D2 parser and grammar implementation.
  \item D2 error reporting, parser recovery, and tests.
  \item D3 malformed programs and written test report.
  \item D5 Enes retrospective.
\end{itemize}

So expect questions on lexical categories, contextual keywords, unit handling, recursive descent, panic-mode recovery, static scoping, lifetimes, and the malformed test outputs.

\section{ANS In One Minute}

\textbf{ANS} means \textbf{Alchemist's Notation Script}. It is a small domain-specific language for chemistry. It lets a user declare elements, compounds, and reactions using chemistry-looking notation:

\begin{lstlisting}[style=ans]
element H { atomic_mass = 1.008, valence = 1, symbol = "H" };
compound water { formula = H2O, name = "Water" };

reaction combustion {
    reactants = CH4 + 2*O2,
    products  = CO2 + 2*H2O,
    check     = balanced,
    energy    = exothermic 890.3 kJ/mol
};
\end{lstlisting}

The domain-specific part is not just the words. It is the first-class syntax for formula literals like \code{H2O}, reaction expressions like \code{CH4 + 2*O2}, and quantity units like \code{kJ/mol}. A generic scripting language could compute the same numbers, but it would hide chemistry inside strings, arrays, or library calls.

\subsection{What Part 1 Does}

\begin{itemize}
  \item Reads an \code{.ans} file.
  \item Lexes it into tokens.
  \item Parses it using a hand-written recursive-descent parser.
  \item Builds an AST.
  \item Prints the AST with \code{--dump-ast}.
  \item Rejects malformed programs with file, line, column, and message.
\end{itemize}

\subsection{What Part 1 Does Not Do}

\begin{itemize}
  \item It does not execute \code{print}, functions, loops, or reactions.
  \item It does not compute molar mass.
  \item It does not type-check values.
  \item It does not validate real chemistry or atom balance.
  \item It does not enforce duplicate/missing declaration fields semantically.
\end{itemize}

If the exam asks, ``Does your parser check that \code{CH4 + 2*O2} is balanced?'', the correct answer is: \textbf{No, Part 1 parses the syntax. Balance checking is a semantic/type-checker responsibility planned for Part 2.}

\section{Sebesta Chapter 1: Design Criteria}

Sebesta's main evaluation criteria are readability, writability, reliability, and cost. Your D1 says ANS deliberately prioritizes readability, writability, and reliability, while sacrificing cost, generality, and some orthogonality.

\begin{tabularx}{\linewidth}{@{}p{3cm}Y Y@{}}
\toprule
\textbf{Criterion} & \textbf{Meaning} & \textbf{ANS answer} \\
\midrule
Readability & How easily code can be understood. & Chemistry is visible: \code{compound water}, \code{H2O}, \code{check = balanced}. \\
Writability & How easily users can express solutions. & Common chemistry tasks are first-class syntax, not manual arrays/strings. \\
Reliability & How well language design prevents errors. & Formula tokens, typed domain parameters, units, and future balance checks catch mistakes early. \\
Cost & Training, implementation, execution, maintenance, reliability cost. & ANS has higher implementation cost because formulas, units, and reactions need custom handling. \\
Generality & How broad the language is. & Sacrificed: ANS is intentionally narrow and chemistry-specific. \\
Orthogonality & Few primitives freely combined in regular ways. & Sacrificed somewhat: ANS has fixed domain forms like \code{reaction} fields. \\
\bottomrule
\end{tabularx}

\subsection{Strong Exam Answer}

\textbf{Question: Why is ANS a good DSL instead of just a generic language?}

ANS is intentionally narrow. It makes chemistry notation part of the language surface. A user writes \code{H2O}, \code{CH4 + 2*O2}, and \code{molar_mass(water)} directly. That improves readability and writability for chemistry tasks. The trade-off is lower generality and higher implementation cost, because the compiler must understand formulas, units, and reaction fields.

\section{Lexical Structure}

The lexer scans characters and produces tokens. In Sebesta terms, lexemes are the actual character sequences, and tokens are categories of lexemes. Example: \code{water} is a lexeme whose token category is \code{IDENT}; \code{H2O} is a lexeme whose token category is \code{FORMULA}.

\subsection{Token Categories}

\begin{tabularx}{\linewidth}{@{}p{3.4cm}p{4cm}Y@{}}
\toprule
\textbf{Category} & \textbf{Examples} & \textbf{Rule} \\
\midrule
Identifier & \code{water}, \code{mass}, \code{compute_yield} & Starts lowercase or underscore: \code{[a-z_][A-Za-z0-9_]*}. \\
Number & \code{20}, \code{1.008}, \code{6.02e23} & Integer/decimal/scientific notation. \\
String & \code{"Water"} & Double quoted, with supported escapes. \\
Char & \code{'H'}, \code{'\n'} & Single quoted character or escape. \\
Formula & \code{H}, \code{O2}, \code{H2O}, \code{CH4} & Starts uppercase, optional lowercase, optional digits, repeated. \\
Unit literal & \code{g/mol}, \code{kg/mol}, \code{kJ/mol} & Slash units recognized before \code{/}. \\
Hard keyword & \code{let}, \code{function}, \code{return}, \code{if} & Identifier-shaped text promoted by keyword table. \\
Contextual word & \code{element}, \code{compound}, \code{reaction} & Lexer leaves as \code{IDENT}; parser interprets by position. \\
Operator & \code{==}, \code{!=}, \code{+=}, \code{+}, \code{/} & Longer operators before shorter ones. \\
Separator & \texttt{\{ \}}, \texttt{( )}, \texttt{[ ]}, \code{,}, \code{;}, \code{.} & Single-character tokens. \\
Comment & \code{// comment} & Skipped by lexer. \\
\bottomrule
\end{tabularx}

\subsection{Formula vs Identifier}

This is one of the most important design choices.

\begin{itemize}
  \item \code{water} starts lowercase, so it is an identifier.
  \item \code{H2O} starts uppercase, so it is a formula literal.
  \item \code{Water} is not an identifier in ANS; it is scanned as formula-like text.
\end{itemize}

This case split keeps chemistry notation visually separate from program names. It also simplifies parsing because the parser can immediately distinguish \code{compound water} from \code{formula = H2O}.

\subsection{Hard Keywords vs Contextual Keywords}

Hard keywords are never valid identifiers:

\begin{center}
\code{let function return for in if elif else print}
\end{center}

Contextual words are special only in grammar positions:

\begin{center}
\code{element compound reaction formula reactants products check balanced energy exothermic endothermic}
\end{center}

Example: in \code{function compute_yield(reaction r, compound c)}, \code{reaction} and \code{compound} are type annotations because they appear before parameter names. Elsewhere, similar words can be ordinary identifiers if the grammar position does not reserve them.

\subsection{Unit Handling}

ANS has two kinds of unit handling:

\begin{itemize}
  \item Slash units \code{g/mol}, \code{kg/mol}, and \code{kJ/mol} are \code{TOK_UNIT_LIT}. The lexer checks them before returning \code{/}.
  \item Single-word units \code{mol}, \code{g}, \code{kg}, and \code{kJ} are scanned as identifiers. The parser treats them as units only right after a number.
\end{itemize}

\begin{lstlisting}[style=ans]
let amount = 5.0 mol;       // NUMBER + IDENT used as unit
let energy = 890.3 kJ/mol;  // NUMBER + UNIT_LIT
\end{lstlisting}

\textbf{Exam phrase:} this avoids a stateful lexer. The lexer stays mostly regular and simple; parser context decides whether an identifier-like word such as \code{mol} is a unit.

\section{Syntax And EBNF}

ANS uses Sebesta-style EBNF. The grammar describes how tokens form valid programs.

\subsection{Top-Level Structure}

A program is a sequence of top-level statements until EOF. Top-level declarations end with semicolons:

\begin{lstlisting}[style=ans]
element H { atomic_mass = 1.008, valence = 1, symbol = "H" };
compound water { formula = H2O, name = "Water" };
function molar_mass(compound c) { return 0.0; };
\end{lstlisting}

Important top-level forms:

\begin{itemize}
  \item Element declarations.
  \item Compound declarations.
  \item Reaction declarations.
  \item Function declarations.
  \item Variable declarations.
  \item Assignments and augmented assignments.
  \item Print and expression statements.
  \item If statements.
\end{itemize}

\subsection{Expression Precedence}

The expression grammar is layered. Higher rows bind more weakly; lower rows bind more tightly.

\begin{tabularx}{\linewidth}{@{}p{4cm}p{4cm}Y@{}}
\toprule
\textbf{Level} & \textbf{Operators/forms} & \textbf{Example} \\
\midrule
Comparison & \code{== != >= <= > <} & \code{mass > 18.0} \\
Additive & \code{+ -} & \code{a + b - c} \\
Multiplicative & \code{* /} & \code{2.0 * 3.0} \\
Unary & \code{-} & \code{-3.14} \\
Postfix & call, index, member & \code{water.name}, \code{f(x)}, \code{a[0]} \\
Primary & literals, identifiers, lists, parentheses & \code{H2O}, \code{[H,O,C]}, \code{(x)} \\
\bottomrule
\end{tabularx}

Example AST evidence from D2:

\begin{lstlisting}[style=ans]
let x = 1.0 + 2.0 * 3.0;
\end{lstlisting}

This parses as \code{1.0 + (2.0 * 3.0)}, because multiplication is deeper in the grammar than addition.

\subsection{Reaction Grammar}

Reaction expressions use a separate grammar from arithmetic expressions.

\begin{lstlisting}[style=ans]
reactants = CH4 + 2*O2
products  = CO2 + 2*H2O
\end{lstlisting}

Inside a reaction:

\begin{itemize}
  \item \code{+} separates species.
  \item \code{*} attaches a coefficient to a species.
  \item Species can be formula literals like \code{CH4} or identifiers like \code{methane}.
  \item ANS requires \code{2*O2}, not \code{2O2}, to avoid mixing coefficients with formula subscripts.
\end{itemize}

\subsection{Ambiguity Answers}

\begin{tabularx}{\linewidth}{@{}p{4cm}Y@{}}
\toprule
\textbf{Potential issue} & \textbf{ANS solution} \\
\midrule
Arithmetic precedence & Separate nonterminals for comparison, additive, multiplicative, unary, postfix, primary. \\
Reaction vs arithmetic & Reactions use \code{reaction_expression}, not the general expression grammar. \\
Coefficient vs subscript & Require explicit \code{2*O2}. \\
Formula vs identifier & Uppercase starts formula; lowercase/underscore starts identifier. \\
Keyword vs identifier & Hard keywords in lexer; contextual words resolved by parser position. \\
Unit vs division & Slash units recognized before \code{/}; single-word units parsed after numbers. \\
Dangling else & Bodies are braced blocks; an \code{else} has a clear nearest \code{if}. \\
\bottomrule
\end{tabularx}

\section{Names, Binding, Scope, Lifetime}

This section is directly tied to your D1 4.6.

\subsection{Legal Identifiers}

An ANS identifier matches:

\begin{center}
\code{[a-z_][A-Za-z0-9_]*}
\end{center}

ANS is case-sensitive. \code{water} and \code{Water} are not the same. In fact, \code{Water} is formula-like, not a normal identifier.

\subsection{Binding Times}

\begin{tabularx}{\linewidth}{@{}p{4.6cm}p{3cm}Y@{}}
\toprule
\textbf{Binding} & \textbf{When} & \textbf{Example} \\
\midrule
Lexeme to token & Lex time & \code{H2O} becomes \code{TOK_FORMULA}. \\
Identifier to declaration & Compile/static analysis time & \code{mass} resolves to a \code{let}. \\
Parameter to declared type & Compile time & \code{compound c}. \\
Formula validity & Part 2 semantic time & Every symbol in \code{H2O} should name an element. \\
Balance check & Part 2 semantic time & \code{check = balanced}. \\
Parameter to argument value & Run time & \code{molar_mass(water)} binds \code{c}. \\
Local variable to storage & Run time & \code{let mass = 0.0}. \\
\bottomrule
\end{tabularx}

\important{Do not say Part 1 already performs semantic formula validation or balance checking.} Say the design plans compile-time validation, but D2 Part 1 implements lexer/parser behavior.

\subsection{Static Scoping}

ANS uses static, lexical scoping. A name resolves by looking at the source text: current block, enclosing block/function, then top-level.

\textbf{Why static?} Chemistry calculations should not change depending on who called the function. The compiler/type checker should be able to resolve names from the program text.

\textbf{What dynamic scoping would break:} a function like \code{molar_mass} could accidentally use a caller's local \code{water}, \code{formula}, or \code{element} binding. Then the same function call could mean different things depending on call chain. That hurts readability, reliability, and compile-time checking.

\subsection{Lifetime Categories}

\begin{tabularx}{\linewidth}{@{}p{4.5cm}p{3.4cm}Y@{}}
\toprule
\textbf{Entity} & \textbf{Lifetime} & \textbf{Reason} \\
\midrule
Top-level elements, compounds, reactions & Static & Program chemistry model exists for whole run. \\
Top-level \code{let} variables & Static & Whole-program lifetime. \\
Function parameters & Stack-dynamic & Each call gets fresh parameter bindings. \\
Local \code{let} variables & Stack-dynamic & Allocated on block/function entry, freed on exit. \\
\code{for} loop variable & Stack-dynamic & Scoped to loop body. \\
Lists and chemistry objects & Implicit heap-dynamic & Runtime manages them; no \code{new}, \code{delete}, or pointers. \\
\bottomrule
\end{tabularx}

\section{D2 Implementation Architecture}

D2 is a C11 project with no external dependencies. Build and run:

\begin{lstlisting}[style=ans]
make
./ans <file.ans>
./ans --dump-ast <file.ans>
make test
\end{lstlisting}

\subsection{File Responsibilities}

\begin{tabularx}{\linewidth}{@{}p{4cm}Y@{}}
\toprule
\textbf{File} & \textbf{Role} \\
\midrule
\code{include/lexer.h} & Token types, \code{Token}, \code{Lexer}, lexer API. \\
\code{src/lexer.c} & Hand-written scanner: comments/whitespace, keywords, strings, chars, numbers, formulas, units, operators. \\
\code{include/parser.h} & \code{Parser} struct and \code{parse_program}. \\
\code{src/parser.c} & Recursive-descent parser and AST construction. \\
\code{include/ast.h} & \code{NodeKind}, \code{AstNode}, AST API. \\
\code{src/ast.c} & AST allocation, child list growth, pretty printing, cleanup. \\
\code{include/error.h} & Error-reporting declarations and globals. \\
\code{src/error.c} & Diagnostic formatting: file, line, column, message. \\
\code{src/main.c} & CLI, file reading, lexer/parser setup, optional AST dump, exit code. \\
\code{Makefile} & Build, debug, test, clean targets. \\
\bottomrule
\end{tabularx}

\subsection{Lexer Pipeline}

The lexer keeps:

\begin{lstlisting}[style=csmall]
typedef struct {
    const char *src;
    int pos;
    int line;
    int col;
} Lexer;
\end{lstlisting}

Each token stores type, source slice, line, column, and numeric value for numbers.

Core lexer logic:

\begin{enumerate}
  \item Skip whitespace and \code{//} comments.
  \item If EOF, return \code{TOK_EOF}.
  \item Probe slash units before ordinary \code{/}.
  \item Scan strings, chars, numbers, identifiers, formulas.
  \item Recognize two-character operators before one-character operators.
  \item Report unknown characters as \code{TOK_ERROR}.
\end{enumerate}

\subsection{Parser Pipeline}

The parser has one-token lookahead:

\begin{lstlisting}[style=csmall]
typedef struct {
    Lexer *lexer;
    Token current;
    Token previous;
} Parser;
\end{lstlisting}

Important helper ideas:

\begin{itemize}
  \item \code{advance}: ask lexer for next token.
  \item \code{check}: test current token type.
  \item \code{match}: consume token if it matches.
  \item \code{expect}: require a token or report an error and recover.
  \item \code{ident_is}: treat contextual words by lexeme comparison.
\end{itemize}

\subsection{Recursive Descent}

Recursive descent means the code structure mirrors the grammar. Example:

\begin{itemize}
  \item \code{parse_expression} calls \code{parse_comparison}.
  \item \code{parse_comparison} calls \code{parse_additive}.
  \item \code{parse_additive} calls \code{parse_multiplicative}.
  \item \code{parse_multiplicative} calls \code{parse_unary}.
  \item \code{parse_unary} calls \code{parse_postfix}.
  \item \code{parse_postfix} calls \code{parse_primary}.
\end{itemize}

This is why the parser is easy to explain in Sebesta terms: one function per grammar level.

\subsection{AST Shape}

Each AST node has:

\begin{itemize}
  \item \code{kind}: node category, such as \code{ND_VAR_DECL}, \code{ND_BINARY}, \code{ND_REACTION_DECL}.
  \item \code{line}: line number for diagnostics/debugging.
  \item \code{text}: name/operator/field/unit string.
  \item \code{num}: numeric payload for number nodes.
  \item \code{children}: dynamically grown list of child pointers.
\end{itemize}

Example AST meaning:

\begin{lstlisting}[style=ans]
let x = 1.0 + 2.0 * 3.0;
\end{lstlisting}

becomes a \code{VarDecl [x]} whose child is \code{Binary [+]}, with right child \code{Binary [*]}. This proves precedence was parsed correctly in that expression context.

\subsection{Error Handling And Recovery}

The error format is:

\begin{center}
\code{<file>:<line>:<col>: error: <message>}
\end{center}

The parser uses panic-mode recovery in \code{expect}: after a missing expected token, it skips until a synchronization token:

\begin{center}
\code{;} \quad \texttt{\}} \quad \code{EOF}
\end{center}

Why this matters: it can report more than one error instead of stopping at the first problem.

\section{D3 Tests}

\subsection{Valid Programs}

\begin{tabularx}{\linewidth}{@{}p{4.4cm}Y@{}}
\toprule
\textbf{Program} & \textbf{What it exercises} \\
\midrule
\code{molar_mass_}\par\code{demo.ans} & Elements, compounds, function definition, \code{for}, \code{if/elif/else}, calls, list literal. \\
\code{combustion_}\par\code{analysis.ans} & Reaction declaration, \code{check = balanced}, energy quantity, typed parameters, nested calls, augmented assignment. \\
\code{multi_reaction_}\par\code{survey.ans} & Multiple reactions, formula and identifier species, loop counters, strings/chars, member access, repeated function calls. \\
\bottomrule
\end{tabularx}

\subsection{Malformed Programs}

\begin{tabularx}{\linewidth}{@{}p{4.5cm}Y@{}}
\toprule
\textbf{Malformed file} & \textbf{Purpose} \\
\midrule
\code{err1_missing_}\par\code{semicolon.ans} & Missing semicolon after element declaration. \\
\code{err2_unknown_}\par\code{field.ans} & Invalid element field \code{protons}. \\
\code{err3_bad_}\par\code{identifier.ans} & Illegal character \code{@}. \\
\code{err4_unmatched_}\par\code{brace.ans} & Missing closing brace inside a function/loop block. \\
\code{err5_unterminated_}\par\code{string.ans} & String reaches newline before closing quote. \\
\bottomrule
\end{tabularx}

\subsection{Observed Verification}

\begin{itemize}
  \item \code{make test} accepts all 8 D2 valid tests.
  \item \code{make test} rejects all 4 D2 invalid tests.
  \item All 3 D3 valid programs parse with exit code 0.
  \item All 5 D3 malformed programs exit with code 1 and line/column diagnostics.
\end{itemize}

\section{Important Caveats}

These are useful because a strong exam answer can admit limitations precisely.

\subsection{Design vs Current Parser: Assignment Target}

\design: D1 says function calls should not be assignment targets, so \code{f(x) = 3.0;} should be invalid.

\actual: the current parser parses assignment from \code{parse_postfix}, and function calls are postfix expressions. Therefore D2 currently accepts:

\begin{lstlisting}[style=ans]
f(x) = 3.0;
\end{lstlisting}

If asked, say: this is an implementation mismatch with the submitted EBNF. It should be fixed by parsing assignment targets with a restricted \code{assignable} parser that allows identifiers, member access, and index access, but not calls.

\subsection{Bare Expression Statements}

\design: general expression statements should use the full expression grammar.

\actual: expressions inside declarations, calls, conditions, lists, and returns use \code{parse_expression}; bare top-level expression statements have a special continuation path. A declaration like:

\begin{lstlisting}[style=ans]
let x = 1.0 + 2.0 * 3.0;
\end{lstlisting}

parses correctly as \code{1 + (2 * 3)}, but a bare statement:

\begin{lstlisting}[style=ans]
1.0 + 2.0 * 3.0;
\end{lstlisting}

is not handled as cleanly by the current special path. The submitted D3 programs avoid this case.

\section{Likely Exam Questions And Model Answers}

\subsection{Language Design}

\textbf{Q: What is ANS?}

ANS is Alchemist's Notation Script, a chemistry DSL. It provides direct syntax for elements, compounds, reactions, formula literals, quantities, and stoichiometric-style function calls. Part 1 implements the lexer, parser, AST dump, and syntax-error reporting.

\textbf{Q: Which Sebesta criteria does ANS prioritize?}

Readability, writability, and reliability. Readability because chemistry notation is visible in the program; writability because common domain tasks are short; reliability because the design gives the compiler places to catch formula, type, unit, and balance mistakes. It sacrifices implementation cost, generality, and some orthogonality.

\textbf{Q: Why is ANS not very orthogonal?}

Because it has fixed domain constructs like \code{element}, \code{compound}, and \code{reaction} with required fields. Those are not freely combinable primitives. That reduces orthogonality but improves domain readability and reliability.

\subsection{Lexical Analysis}

\textbf{Q: How does the lexer distinguish formulas from identifiers?}

By the first character. Identifiers start lowercase or underscore. Formula literals start uppercase. So \code{water} is an identifier and \code{H2O} is a formula literal.

\textbf{Q: Why not make \code{mol} a hard keyword?}

Because \code{mol} may be useful as an ordinary name in some contexts. The lexer scans it as \code{IDENT}; the parser treats it as a unit only when it immediately follows a number.

\textbf{Q: Why are \code{g/mol} and \code{kJ/mol} special?}

They contain \code{/}, which is also division. The lexer probes slash units before returning \code{/}. This is longest-match behavior and prevents \code{kJ/mol} from becoming \code{kJ}, \code{/}, \code{mol}.

\subsection{Syntax And Parsing}

\textbf{Q: Why recursive descent?}

The grammar is written in a recursive-descent-friendly style, with separate functions for grammar levels. It is simple to explain, directly tied to the EBNF, and appropriate for Part 1's hand-written C parser.

\textbf{Q: What problem does left recursion cause?}

In recursive descent, a left-recursive rule calls itself before consuming input, causing infinite recursion or stack overflow. ANS avoids left recursion by using EBNF repetitions and loop-based parsing.

\textbf{Q: How do you handle dangling else?}

All ANS branch bodies are braced blocks. Because each \code{if}, \code{elif}, and \code{else} body is clearly delimited, there is no ambiguous unbraced statement for an \code{else} to attach to.

\textbf{Q: Why require \code{2*O2}?}

Without \code{*}, \code{2O2} would blur stoichiometric coefficients and formula notation. Requiring \code{2*O2} makes the coefficient syntactically explicit and keeps formula literals uppercase-led.

\subsection{Scope And Binding}

\textbf{Q: Static or dynamic scoping?}

ANS uses static scoping. Names resolve from the source text, not the call chain. This supports readability and future compile-time checking.

\textbf{Q: What would break with dynamic scoping?}

A function could refer to a caller's local variable accidentally. For example, \code{molar_mass} could use a caller's \code{formula} or \code{water} binding instead of the intended one. That makes chemistry calculations depend on call history and prevents reliable compile-time checking.

\textbf{Q: What is the lifetime of local variables?}

Function parameters and local \code{let} variables are stack-dynamic. They are allocated per call or block execution and disappear when the call/block ends. Top-level declarations have static lifetime.

\subsection{Implementation}

\textbf{Q: What does \code{--dump-ast} show?}

It prints the tree structure built by the parser. It shows declarations, fields, expressions, reaction terms, and operator hierarchy. It is evidence that parsing works for that input, but it is not a proof of semantic correctness.

\textbf{Q: What is panic-mode recovery?}

When \code{expect} does not see the token it needs, it reports an error and skips tokens until a synchronization point such as semicolon, right brace, or EOF. That lets the parser continue and report additional errors.

\textbf{Q: What is the role of header files?}

Headers define module interfaces. \code{lexer.h} exposes token and lexer types; \code{parser.h} exposes the parser API; \code{ast.h} exposes AST structures/functions; \code{error.h} exposes diagnostics. Implementations live in the corresponding \code{.c} files.

\section{Three Handwritten A4 Sheet Summary}

\subsection*{Sheet 1: Language And Sebesta}

\begin{itemize}
  \item ANS = chemistry DSL: elements, compounds, reactions, formulas, units.
  \item Part 1 = lexer + parser + AST + diagnostics, not interpreter/type checker.
  \item Prioritizes readability, writability, reliability.
  \item Sacrifices cost, generality, orthogonality.
  \item DSL justification: \code{H2O}, \code{CH4 + 2*O2}, \code{check = balanced} are direct syntax.
  \item Syntax vs semantics: syntax is form; semantics is meaning. Part 1 mainly checks syntax.
  \item Lexeme vs token: \code{water} lexeme, \code{IDENT} token; \code{H2O} lexeme, \code{FORMULA} token.
\end{itemize}

\subsection*{Sheet 2: Lexer, Grammar, Parser}

\begin{itemize}
  \item Identifier regex: \code{[a-z_][A-Za-z0-9_]*}.
  \item Formula regex idea: uppercase element groups with optional lowercase and digits.
  \item Hard keywords: \code{let function return for in if elif else print}.
  \item Contextual words: \code{element compound reaction balanced exothermic}.
  \item Slash units: \code{g/mol kg/mol kJ/mol} before \code{/}.
  \item Single units: \code{mol g kg kJ} are identifiers, parsed as units after numbers.
  \item Precedence: comparison < additive < multiplicative < unary < postfix < primary.
  \item Reaction grammar separate: \code{+} species separator, \code{*} coefficient marker.
  \item Recursive descent = one function per grammar rule/level.
  \item Avoid left recursion for recursive descent.
  \item Panic recovery skips to \code{;}, \texttt{\}}, EOF.
\end{itemize}

\subsection*{Sheet 3: Scope, Code, Tests, Caveats}

\begin{itemize}
  \item Static scoping: resolve by source text. Dynamic scoping would depend on caller chain.
  \item Bindings: lexeme-token at lex time; declarations/types/formula checks at compile/semantic time; values/storage at run time.
  \item Lifetimes: top-level static; params/local lets stack-dynamic; objects/lists implicit heap-dynamic.
  \item D2 files: \code{lexer.c}, \code{parser.c}, \code{ast.c}, \code{error.c}, \code{main.c}.
  \item CLI: \code{./ans file.ans}, \code{./ans --dump-ast file.ans}, \code{make test}.
  \item D3 valid: molar mass demo, combustion analysis, multi-reaction survey.
  \item D3 malformed: missing semicolon, unknown field, illegal \code{@}, unmatched brace, unterminated string.
  \item Caveat: D1 says calls not assignable; current parser accepts \code{f(x)=3}.
  \item Caveat: semantic balance/type checks are Part 2 design, not D2 Part 1 behavior.
\end{itemize}

\section{Final Defense Script}

If you freeze in the exam, use this structure:

\begin{enumerate}
  \item State whether the question is about design, grammar, lexer, parser, or future semantics.
  \item Give the ANS example: \code{H2O}, \code{water}, \code{CH4 + 2*O2}, \code{5.0 mol}.
  \item Name the Sebesta concept: readability, token, BNF/EBNF, static scoping, lifetime, binding.
  \item Say what D2 actually implements.
  \item If there is a limitation, admit it precisely and say how you would fix it.
\end{enumerate}

Example final answer style:

\begin{quote}
In the design, ANS wants \code{f(x) = 3} to be illegal because assignment targets are supposed to be identifiers with optional member/index access. In the current D2 parser, assignment starts from \code{parse_postfix}, so a call can accidentally become an assignment target. That is a parser/grammar mismatch, not a language-design decision. I would fix it by implementing a separate \code{parse_assignable} function that excludes call postfixes.
\end{quote}

\end{document}
